\documentclass{article}
\usepackage{lmodern}
\begin{document}

\section{The problem}

A research console can produce a confident paragraph and a citation that does not exist in
the material the researcher actually registered. The damage is not the prose: it is that a
reader has no cheap way to see which sentences the stored evidence supports.

A Claim-Evidence check ships in this build, and it is implemented and independently
reviewed here: the engine is \texttt{app/claim\_check.py} and the interface is
\texttt{frontend/claim\_check\_ui.js}. Numbers in a draft are reconciled against a stored results
table by ordinary code with no model involved, so the same two arguments always give the
same answer and no cost is incurred. Statements about a stored source are put to a model,
which returns the passage it says supports or contradicts the statement, and app code
checks that passage verbatim against the excerpt actually sent, so a quote the excerpt does
not contain cannot be shown as evidence. No model call happens without an exact preview of
what will be sent, where, on which model and at what cost, and the researcher's consent.
Neither step decides whether an argument is good. What comes back is a secretary's note: a
hedged, first person suggestion written in the researcher's language, which the researcher
decides on, and never a ruling that a sentence is true. When no quote survives the verbatim
check, the app says so in its own words instead of showing the model's note.

\section{What is measured here}

The console runs on one local SQLite database and one local interface. A local 4B
instruction model can serve as a zero-cost arm: the app starts the installed server itself,
verifies its identity, and stops it on quit. Nothing about that arm is claimed as a quality
win. The claim check also runs in the cloud, on Baseten Muse Spark Open, and that path is
live: the first live call, job aichk-d6f42801700a, returned a supported verdict with one
quote verified verbatim in 3.4 s, for a calculated 0.004114 USD on 1,105 input and 643
output tokens. It also exposed two defects, since fixed and independently reviewed: an
internal source ID appeared where the source title belonged, and Korean rendered the word
for an assertion as the word for a billing claim. A later benchmark of five alternating
calls per arm showed a median of 2.55 s at reasoning effort low against 2.11 s at none; the
0.45 s gap sits inside run-to-run spread and caching favored the cheaper arm, so effort
stays low and the difference is not reported as a speed result. All ten calls were
supported with one verified quote each, and none leaked an ID or used the billing word.

\subsection{Current round 3}

AI 1차 채점 → 사람 2차 검수(출처 공개 뒤)

Ten questions across 3 FX, 4 DOB NOW, 3 BIS yielded 56/60 retained answers,
all AI-scored. Human round-3 review is pending.

\begin{center}
\begin{tabular}{lrr}
Provider & Basic success / planned & Harness success / planned \\
Baseten & 9/10 & 9/10 \\
Local Qwen & 4/10 & 6/10 \\
Modal Shared & 7/10 & 7/10 \\
\end{tabular}
\end{center}

Calculated round spend: Baseten 0.13738895 USD; Modal 0.21011545 USD;
local API 0 USD. Baseten cross-round total is 0.20959970 USD.
Modal retains an unresolved 3.958272 USD reservation, giving conservative
exposure of 4.16838745 USD. Neither approved cap was exceeded; invoices
are not verified. Two Modal jobs failed and one partial answer was discarded.
Later Modal jobs continued contrary to the unknown-charge stop rule.
Correlated datasets, provider profiles and one repetition limit the result:
this is not established harness superiority. See the round-3 post-run
preregistration and \texttt{ai-proxy-scores.json}.

\subsection{Historical round 2}

A basic-versus-harness research comparison ran on 2026-11-01: Baseten
\texttt{meta-models/muse-spark-open-1026} and local Qwen3-4B, n=1 per case-arm.
All 9 Baseten calls completed, at a calculated cost of about 0.0722 USD.
Scoring: AI first-pass scoring, human second-pass review (after unblinding).
GPT-6.1-Sol medium blind-scored all 36 claims and 80 citation occurrences.
The author-reviewer agreed with those scores, changing no marks after unblinding;
the second pass was not blinded or independent.

Overall ratings (1 to 5), from the Post-run addendum in
\texttt{HARNESS\_COMPARISON\_PREREG\_EN.md} and \texttt{ai-proxy-scores.json}:

\begin{center}
\begin{tabular}{lrrrr}
Case & Baseten basic & Baseten harness & Local basic & Local harness \\
FX-COST-1 & 4 & 4 & 4 & 4 \\
FX-APPROVAL-1 & 4 & 5 & 4 & 3 \\
NYC-DATE-1 & 3 & 4 & 2 & 3 \\
\end{tabular}
\end{center}

On Baseten, the harness rated higher on 2 of 3 cases and tied on 1.
Local Qwen results were mixed: higher on 1, lower on 1, tied on 1.
All four FX-COST-1 answers omitted ``not actual final cost''.
This is an indication, not an established harness effect: n=1, AI first pass,
unblinded second pass, thresholds never set. Preregistered aggregates remain NOT_SCORED.

The 4.6 to 6.4 second band below is for short questions; a claim check that sends source
excerpts takes longer on the same local path, 8.0 to 11.1 seconds, still at zero cost.

\section{Claims under verification}

Each sentence below is one the check is run against in the demo, judged against this
project's own stored notes and results table. The first three are the qualitative case; the rest are
numbers already recorded in the project's results table in Literature Review, stored as
measurements.csv.

\begin{enumerate}
  \item Local mode makes no API call, so a rehearsal costs nothing in API credits.
  \item The app stops the local model server it started when the app quits.
  \item Grounding checks confirm that a sentence is true.
  \item In rehearsal the local model answered in 4.6 to 6.4 seconds per question.
  \item The app declares an 8192-token context bound, the chat template overhead observed
        at launch is 8 tokens, and the installed weights file is 2,497,281,120 bytes.
  \item The full Python suite reached 2,224 passed and 1 xfailed.
  \item An earlier count recorded 2,300 passed.
\end{enumerate}

\section{What this draft does not claim}

It does not claim that the cloud route is unproven: one live claim check ran there, on
Baseten Muse Spark Open, and is reported above with its job id and its calculated cost.
It does not claim live cloud quality from that single call either, because one supported
verdict is not a measurement of quality, and every cost stated here is calculated rather
than read off an invoice. The comparison above is not an established harness effect.
Self-hosted Spark Open and Glimmer benchmarks remain NOT_RUN; SAM mask raster decoding
is NOT_IMPLEMENTED. No submission has been filed and no public repository created.
The local route was verified live in the packaged build on
2026-10-31. It does not claim model quality, a comparison win, or that any statement here
is free of error. The check reports what it can see and states the rest as unverified.

\end{document}
